% Part:sets-functions-relations
% Chapter: size-of-sets
% Section: reduction

\documentclass[../../../include/open-logic-section]{subfiles}

\begin{document}

\olfileid{sfr}{siz}{red}

\olsection{తగ్గింపు}

\begin{editorial}
  ఈ విభాగం, \olref[nen]{sec}లోని ఫలితాలకు సరిపోలేలా తగ్గింపు ద్వారా
  లెక్కించలేనితనాన్ని నిరూపిస్తుంది. \olref[nen-alt]{sec}లోని ఫలితాలకు
  సరిపోలే, కొంచెం సంక్షిప్తమైన ప్రత్యామ్నాయ రూపాన్ని
  \olref[red-alt]{sec}లో ఇచ్చాం.
\end{editorial}

వికర్ణీకరణ వాదనతో $\Pow{\PosInt}$ \tetoken{లెక్కించలేని}{!!{nonenumerable}} అని చూపాం.
$0$, $1$ల అనంత క్రమాలన్నిటి సమితి $\Bin^\omega$ కూడా
\tetoken{లెక్కించలేని}{!!{nonenumerable}} అని ఇప్పటికే నిరూపించాం. $\Pow{\PosInt}$
\tetoken{లెక్కించలేని}{!!{nonenumerable}} అని నిరూపించడానికి మరో మార్గం ఇదిగో:
\emph{$\Pow{\PosInt}$ \tetoken{లెక్కించదగిన}{!!{enumerable}} అయితే $\Bin^\omega$ కూడా
  \tetoken{లెక్కించదగిన}{!!{enumerable}} అవుతుందని చూపడం}. $\Bin^\omega$ \tetoken{లెక్కించదగిన}{!!{enumerable}}
కాదని మనకు తెలుసు కనుక, $\Pow{\PosInt}$ కూడా అలా ఉండజాలదు.
దీనినే ఒక సమస్యను మరొక సమస్యకు \emph{తగ్గించడం} అంటారు---ఇక్కడ
$\Bin^\omega$ను లెక్కించే సమస్యను $\Pow{\PosInt}$ను లెక్కించే
సమస్యకు తగ్గిస్తాం. రెండవ సమస్యకు పరిష్కారం---$\Pow{\PosInt}$ యొక్క
లెక్కింపు---మొదటి సమస్యకు పరిష్కారాన్ని---$\Bin^\omega$ యొక్క
లెక్కింపును---ఇస్తుంది.

$B$ సమితిని లెక్కించే సమస్యను $A$ సమితిని లెక్కించే సమస్యకు ఎలా
తగ్గిస్తాం? $A$ యొక్క లెక్కింపును $B$ యొక్క లెక్కింపుగా మార్చే ఒక
విధానాన్ని ఇస్తాం. అందుకు అత్యంత సులభమైన మార్గం \tetoken{ఒక సంగ్రస్త}{!!a{surjective}}
ప్రమేయం $f\colon A \to B$ను నిర్వచించడం. $x_1$, $x_2$, \dots{}
$A$ను లెక్కిస్తే, $f(x_1)$, $f(x_2)$, \dots{} $B$ను లెక్కిస్తాయి.
మన సందర్భంలో, $f\colon \Pow{\PosInt} \to \Bin^\omega$ అనే సంగ్రస్త
ప్రమేయం కోసం చూస్తున్నాం.

\begin{prob}
$g\colon B \to A$ అనే \tetoken{ఒకటి-ఒకటి}{!!{injective}} ప్రమేయం ఉండి, $B$~\tetoken{లెక్కించలేని}{!!{nonenumerable}}
అయితే $A$ కూడా లెక్కించలేనిది అని చూపండి. $A$ యొక్క లెక్కింపును
$B$ యొక్క లెక్కింపుగా మార్చడానికి $g$ను ఎలా వాడవచ్చో చూపి దీన్ని
నిరూపించండి.
\end{prob}

\begin{proof}[తగ్గింపు ద్వారా {\olref[nen]{thm:nonenum-pownat}} నిరూపణ]
$\Pow{\PosInt}$ \tetoken{లెక్కించదగిన}{!!{enumerable}} అనుకుందాం; అందువల్ల దానికి
$Z_{1}$, $Z_{2}$, $Z_{3}$, \dots అనే లెక్కింపు ఉందనుకుందాం.

$f(Z)$ను $s$ అనే క్రమంగా తీసుకుని, $n \in Z$ అయినప్పుడు మాత్రమే
$s(n) = 1$, లేకపోతే $s(n) = 0$ అయ్యేలా
$f \colon \Pow{\PosInt} \to \Bin^\omega$ను నిర్వచిద్దాం. ఇది
స్పష్టంగా ఒక ప్రమేయాన్ని నిర్వచిస్తుంది: $Z \subseteq \PosInt$
అయినప్పుడల్లా, ఏ $n \in \PosInt$ అయినా $Z$ యొక్క \tetoken{ఒక మూలకం}{!!a{element}}
అవుతుంది లేదా కాదు. ఉదాహరణకు, సరి ధన పూర్ణాంకాల సమితి
$2\PosInt = \{2, 4, 6, \dots\}$ను $010101\dots$ క్రమానికి,
ఖాళీ సమితిని $0000\dots$కు, $\PosInt$నే $1111\dots$కు పంపుతుంది.
\sourcecorrection{OLSIZ-004}{ఆంగ్ల మూలం $f(Z)$ ఫలితాన్ని $s_k$గా
పిలిచి, $k$ను ఎక్కడా బంధించలేదు. ఇక్కడ అదే ఫలిత క్రమానికి $s$ అనే
ఒక్క పేరును నిలకడగా వాడి నిర్వచనాన్ని సరిచేశాం.}

ఇది \tetoken{సంగ్రస్త}{!!{surjective}} కూడా: $0$, $1$ల ప్రతి క్రమానికి ఏదో ఒక ధన
పూర్ణాంకాల సమితి సరిపోతుంది; ఆ క్రమంలో $1$లు ఉన్న స్థానాలకు సరిపడే
పూర్ణాంకాలనే ఆ సమితి మూలకాలుగా తీసుకుంటాం. మరింత కచ్చితంగా,
$s \in \Bin^\omega$ అనుకుందాం. $Z \subseteq \PosInt$ను ఇలా
నిర్వచిద్దాం:
\[
Z = \Setabs{n \in \PosInt}{s(n) = 1}
\]
అప్పుడు $f$ నిర్వచనాన్ని పరిశీలిస్తే $f(Z) = s$ అని తెలుస్తుంది.

ఇప్పుడు ఈ జాబితాను పరిశీలించండి:
\[
f(Z_1), f(Z_2), f(Z_3), \dots
\]
$f$ \tetoken{సంగ్రస్త}{!!{surjective}} కనుక, $\Bin^\omega$ యొక్క ప్రతి మూలకం ఏదో ఒక
ఆర్గ్యుమెంటుకు $f$ విలువగా రావాలి; అందువల్ల ఈ జాబితాలో రావాలి.
కాబట్టి ఈ జాబితా $\Bin^\omega$ అంతటినీ లెక్కిస్తుంది.

అందుచేత $\Pow{\PosInt}$ \tetoken{లెక్కించదగిన}{!!{enumerable}} అయితే $\Bin^\omega$ కూడా
\tetoken{లెక్కించదగిన}{!!{enumerable}} అవుతుంది. కానీ $\Bin^\omega$ \tetoken{లెక్కించలేని}{!!{nonenumerable}}
(\olref[nen]{thm:nonenum-bin-omega}). కాబట్టి $\Pow{\PosInt}$
\tetoken{లెక్కించలేని}{!!{nonenumerable}}.
\end{proof}

\begin{explain}
తగ్గింపు ఏ దిశలో సాగుతుందో విషయంలో గందరగోళం పడటం సులభం.
ఉదాహరణకు, \tetoken{ఒక సంగ్రస్త}{!!a{surjective}} ప్రమేయం $g \colon \Bin^\omega \to B$
ఉండటం వల్ల $B$ \tetoken{లెక్కించలేని}{!!{nonenumerable}} అని \emph{నిర్ధారణ కాదు}.
($g(s) = s(1)$గా నిర్వచించిన $g \colon \Bin^\omega \to \Bin$ను
పరిశీలించండి; ఇది $0$, $1$ల క్రమాన్ని దాని మొదటి \tetoken{మూలకానికి}{!!{element}}
పంపుతుంది. కొన్ని క్రమాలు $0$తో, మరికొన్ని $1$తో మొదలవుతాయి కనుక
ఇది \tetoken{సంగ్రస్త}{!!{surjective}}. కానీ $\Bin$ పరిమిత సమితి.) అంతేకాక, $f$ తప్పక
\tetoken{సంగ్రస్త}{!!{surjective}} కావాలి; లేకపోతే వాదన పనిచేయదు: $f(x_1)$, $f(x_2)$,
\dots{}లో $B$ యొక్క \tetoken{మూలకాలు}{!!{element}s} అన్నీ ఉంటాయని హామీ ఉండదు. ఉదాహరణకు,
\[
h(n) = \underbrace{000\dots0}_{\text{$n$ సున్నాలు}}111\dots
\]
$h\colon \PosInt \to \Bin^\omega$ అనే ప్రమేయాన్ని నిర్వచిస్తుంది;
అయినా $\PosInt$ \tetoken{లెక్కించదగిన}{!!{enumerable}}.
\sourcecorrection{OLSIZ-005}{ఆంగ్ల మూలం $h(n)$ను కేవలం $n$ సున్నాల
పరిమిత పదబంధంగా ఇచ్చి, దానిని $\Bin^\omega$లోని అనంత క్రమంగా
టైప్ చేసింది. ఇక్కడ ఆ $n$ సున్నాల తరువాత అనంతంగా $1$లు చేర్చి
నిజమైన ద్విమితీయ అనంత క్రమంగా సరిచేశాం; ఉదాహరణ ఉద్దేశం మారలేదు.}
\end{explain}

\begin{prob}
ధన పూర్ణాంకాల జతల \emph{సమితులన్నిటి} సమితి \tetoken{లెక్కించలేని}{!!{nonenumerable}} అని
తగ్గింపు వాదనతో చూపండి.
\end{prob}

\begin{prob}\label{sfr:siz:red:prob:nat-nat}
  $f\colon \Nat \to \Nat$ రూపంలోని ప్రమేయాలన్నిటి సమితి $X$
  \tetoken{లెక్కించలేని}{!!{nonenumerable}} అని తగ్గింపు వాదనతో చూపండి. (సూచన: $X$ నుంచి
  $\Bin^\omega$పైకి ఒక సంగ్రస్త ప్రమేయాన్ని ఇవ్వండి.)
\end{prob}

\begin{prob}
సహజ సంఖ్యల అనంత క్రమాల సమితి $\Nat^\omega$ \tetoken{లెక్కించలేని}{!!{nonenumerable}} అని
తగ్గింపు వాదనతో చూపండి.
\end{prob}

\begin{prob}
ధన పూర్ణాంకాల సమితి నుంచి $\{0\}$ సమితికి ప్రమేయాల సమితి $P$,
ధన పూర్ణాంకాల సమితి నుంచి $\{0\}$కు \emph{పాక్షిక} ప్రమేయాల సమితి
$Q$ అనుకుందాం. $P$~\tetoken{లెక్కించదగిన}{!!{enumerable}}, $Q$ అలా కాదని చూపండి. (సూచన:
$\Bin^\omega$ను లెక్కించే సమస్యను $Q$ను లెక్కించే సమస్యకు తగ్గించండి.)
\end{prob}

\begin{prob}
ధన పూర్ణాంకాల సమితి నుంచి \{0,1\} సమితిపైకి \tetoken{సంగ్రస్త}{!!{surjective}}
ప్రమేయాలన్నిటి సమితి $S$ అనుకుందాం; అంటే, \tetoken{సంగ్రస్త}{!!{surjective}}
$f\colon \PosInt \to \Bin$
ప్రమేయాలన్నిటితో $S$ ఏర్పడుతుంది. $S$ \tetoken{లెక్కించలేని}{!!{nonenumerable}} అని చూపండి.
\end{prob}

\begin{prob}
వాస్తవ సంఖ్యల సమితి $\Real$ \tetoken{లెక్కించలేని}{!!{nonenumerable}} అని చూపండి.
\end{prob}

\end{document}
